% Discussion, limitations, data and code. Numbers from sections/numbers-lttv.tex.
% Draft, 2026-10-03.

\section{What the case shows}
\label{sec:discussion}

The report's technical claim held in one place and failed in another. At the pre-stated reading, distributors lost more revenue than the model predicted, though the interval is wide and the entry-level service drew far fewer subscribers than the model assumed, so how much of that loss the CRTC's decisions caused isn't settled here. The chain then broke at the transfer to Canadian channels, whose revenue stayed at or above what the report expected with no reform at all.

The jobs figure put to Parliament rests on that chain. Of the \lttvJobsTotal{} FTEs, \lttvFTESpecProd{} (\lttvFTESpecProdShare{}) are in specialty and pay services and independent production, and depend in the model on Canadian channels' revenue and programming spending; the other \lttvFTEBDU{} are distributors' own \citep[Tables~22--23, pp.~91--92]{nordicity2015television}. The staff counts show employment moving against that chain: channels shed more staff than the forecast's direct loss although their revenue held, and distributors didn't shed staff although their revenue fell. Channel jobs did fall, then, but not by the route the model describes, and the distributor jobs the model tied to the revenue loss didn't go. Spin-off employment has no observable counterpart, and production employment remains to be checked.

The public argument went further than the technical claim. The testimony presented an economy-wide estimate, more than half of it spin-off employment, as media jobs that would be lost, as the direct result of the policy, and it quoted a share of programming spending that the report's own figures don't support. On the way from report to committee, a conditional projection with stated assumptions became a point prediction without them. It's a particular case of the point \textcite[4]{manski2011certitude} makes about policy analysis generally, where point predictions are common and statements of uncertainty rare.

The report deserves credit for what made this check possible. It stated its assumptions link by link, published annual paths for both scenarios, cited the range of outside estimates it chose among, and gave the figures behind its headline. A forecast of a single number would have left nothing to check but the number. The useful lesson for those who commission or hear such forecasts is in that contrast. Intermediate predictions that public data can test, an uptake share or a payment flow in the second year, let a forecast be held to account before its end date, and a statement of how wrong the baseline could be shows in advance how much the headline depends on it.

\subsection{Limits of the check}

The paper checks one forecast, and nothing here shows that stakeholder forecasts in general are biased; \textcite{simpson2014overestimate} explains why even a run of overestimates needn't show that. The verdict on channel revenue rests on how wrong the report's technology-only baseline is taken to be. I model that error as a random walk calibrated in several ways, the simplest model that lets the uncertainty be stated, and a structural break in viewing could exceed it. Entry-level uptake is known only for 2016, closures of independent channels can't be separated from moves to exemption, the aggregates show that the pass-through didn't appear without showing why, and the staff counts are a different measure from the report's modelled employment. Canadian programming expenditure moved too much from year to year, before the reforms, to test the forecast on that measure.

None of these limits touches the two findings that need no counterfactual. Payments to Canadian channels rose while the report's own rule says they should have fallen, and the testimony stated more than the report did.

\section*{Data and code}

Every number in this paper can be reproduced from the repository at \url{https://github.com/BrettRey/prediction-checker}, which holds a script that downloads every public source used, the frozen prediction record, the scripts that read the report's figures and check them against its text, the design analysis, and the outcome analysis, with the CRTC series used (Open Government Licence~-- Canada). The analysis plan and each choice made after the outcome data was opened are recorded with dates in the repository. The case was first suggested in a list of candidate cases produced in a ChatGPT research session; every claim taken from that list was checked against primary sources before use.
